%%%%%%%%%%%%%%%%%%%%%%%%%%%%%%%%%%%%%%%%%%%%%%%%%%%%%%%%%%%%
% 4.5 MANIFOLDS
% The Composition of Advanced Mathematics
%%%%%%%%%%%%%%%%%%%%%%%%%%%%%%%%%%%%%%%%%%%%%%%%%%%%%%%%%%%%

\section{Manifolds}
\label{sec:manifolds}

\prerequisite{
Point-Set Topology, Algebraic Topology, Differential Topology,
Geometric Topology, Linear Algebra, and Multivariable Calculus.
}

\learningobjectives{
By the end of this section, the reader should be able to:
\begin{itemize}
    \item define topological, smooth, and differentiable manifolds;
    \item explain local Euclidean structure;
    \item distinguish local dimension from ambient dimension;
    \item define coordinate charts and atlases;
    \item explain transition maps and compatibility;
    \item describe smooth structures and maximal atlases;
    \item define manifolds with boundary;
    \item identify interiors and boundaries of manifolds;
    \item describe submanifolds;
    \item construct product manifolds;
    \item describe quotient constructions that produce manifolds;
    \item explain tangent spaces and tangent bundles;
    \item define cotangent spaces and cotangent bundles;
    \item explain vector fields and differential forms;
    \item define smooth maps and diffeomorphisms;
    \item explain immersions, submersions, and embeddings;
    \item describe orientability;
    \item explain partitions of unity;
    \item state the Whitney Embedding Theorem;
    \item describe Riemannian metrics on smooth manifolds;
    \item explain Lie groups as manifolds with algebraic structure;
    \item introduce fiber bundles and sections;
    \item explain vector bundles;
    \item describe covering spaces as locally trivial constructions;
    \item explain characteristic classes conceptually;
    \item describe compactness, connectedness, and completeness;
    \item distinguish topological, PL, and smooth manifold categories;
    \item recognize important examples of manifolds.
\end{itemize}
}

%%%%%%%%%%%%%%%%%%%%%%%%%%%%%%%%%%%%%%%%%%%%%%%%%%%%%%%%%%%%
\subsection{Why Manifolds?}
\label{subsec:why-manifolds}
%%%%%%%%%%%%%%%%%%%%%%%%%%%%%%%%%%%%%%%%%%%%%%%%%%%%%%%%%%%%

Many mathematical spaces are globally complicated but locally simple.

The surface of the Earth is curved globally, yet a sufficiently small region
looks approximately flat.

A circle is globally closed, but near any point it looks like an interval.

A sphere is globally different from the plane, but each sufficiently small
piece looks like an open subset of the plane.

This principle motivates manifolds:

\[
\boxed{
\text{globally complicated}
\qquad
\text{but}
\qquad
\text{locally Euclidean}.
}
\]

\begin{conceptbox}
A manifold is a space that locally behaves like Euclidean space even when
its global topology or geometry is very different from Euclidean space.
\end{conceptbox}

%%%%%%%%%%%%%%%%%%%%%%%%%%%%%%%%%%%%%%%%%%%%%%%%%%%%%%%%%%%%
\subsection{Topological Manifolds}
\label{subsec:topological-manifolds}
%%%%%%%%%%%%%%%%%%%%%%%%%%%%%%%%%%%%%%%%%%%%%%%%%%%%%%%%%%%%

\begin{definitionbox}{Topological \(n\)-Manifold}
A \emph{topological \(n\)-manifold} is a topological space \(M\) satisfying:

\begin{enumerate}
    \item \(M\) is Hausdorff;
    \item \(M\) is second countable;
    \item every point \(p\in M\) has a neighborhood homeomorphic to an open
    subset of \(\R^n\).
\end{enumerate}
\end{definitionbox}

The integer

\[
\boxed{
n
}
\]

is called the dimension of the manifold.

%%%%%%%%%%%%%%%%%%%%%%%%%%%%%%%%%%%%%%%%%%%%%%%%%%%%%%%%%%%%
\subsection{Local Euclidean Structure}
\label{subsec:local-euclidean-structure}
%%%%%%%%%%%%%%%%%%%%%%%%%%%%%%%%%%%%%%%%%%%%%%%%%%%%%%%%%%%%

For every point

\[
p\in M,
\]

there exists an open neighborhood

\[
U\subseteq M
\]

and a homeomorphism

\[
\boxed{
\varphi:U\to\varphi(U)\subseteq\R^n,
}
\]

where

\[
\varphi(U)
\]

is open in \(\R^n\).

Thus locally,

\[
\boxed{
M\sim\R^n.
}
\]

This does not mean that \(M\) is globally homeomorphic to \(\R^n\).

%%%%%%%%%%%%%%%%%%%%%%%%%%%%%%%%%%%%%%%%%%%%%%%%%%%%%%%%%%%%
\subsection{Why Hausdorff?}
\label{subsec:why-hausdorff-manifolds}
%%%%%%%%%%%%%%%%%%%%%%%%%%%%%%%%%%%%%%%%%%%%%%%%%%%%%%%%%%%%

The Hausdorff condition ensures that distinct points can be separated by
disjoint neighborhoods.

If

\[
p\neq q,
\]

there exist open sets \(U\) and \(V\) such that

\[
p\in U,
\qquad
q\in V,
\qquad
U\cap V=\varnothing.
\]

This prevents many pathological spaces from being considered manifolds.

%%%%%%%%%%%%%%%%%%%%%%%%%%%%%%%%%%%%%%%%%%%%%%%%%%%%%%%%%%%%
\subsection{Why Second Countable?}
\label{subsec:why-second-countable-manifolds}
%%%%%%%%%%%%%%%%%%%%%%%%%%%%%%%%%%%%%%%%%%%%%%%%%%%%%%%%%%%%

A space is second countable if its topology has a countable basis.

This assumption ensures useful global properties.

For manifolds it contributes to results involving

\[
\boxed{
\text{paracompactness},
\quad
\text{partitions of unity},
\quad
\text{countable atlases}.
}
\]

Without such conditions, locally Euclidean spaces can behave pathologically.

%%%%%%%%%%%%%%%%%%%%%%%%%%%%%%%%%%%%%%%%%%%%%%%%%%%%%%%%%%%%
\subsection{Dimension}
\label{subsec:manifold-dimension}
%%%%%%%%%%%%%%%%%%%%%%%%%%%%%%%%%%%%%%%%%%%%%%%%%%%%%%%%%%%%

The dimension of a manifold is determined by the Euclidean space appearing
in its local models.

For example,

\[
S^1
\]

is \(1\)-dimensional,

\[
S^2
\]

is \(2\)-dimensional,

and

\[
S^n
\]

is \(n\)-dimensional.

Although

\[
S^n\subseteq\R^{n+1},
\]

its manifold dimension is

\[
\boxed{
n,
}
\]

not \(n+1\).

%%%%%%%%%%%%%%%%%%%%%%%%%%%%%%%%%%%%%%%%%%%%%%%%%%%%%%%%%%%%
\subsection{Intrinsic Versus Ambient Dimension}
\label{subsec:intrinsic-ambient-dimension}
%%%%%%%%%%%%%%%%%%%%%%%%%%%%%%%%%%%%%%%%%%%%%%%%%%%%%%%%%%%%

The dimension of a manifold is intrinsic.

For example, a curve may be embedded as

\[
S^1\subseteq\R^2
\]

or as

\[
S^1\subseteq\R^3.
\]

In either case,

\[
\boxed{
\dim S^1=1.
}
\]

The ambient dimension depends on the chosen embedding.

The manifold dimension does not.

%%%%%%%%%%%%%%%%%%%%%%%%%%%%%%%%%%%%%%%%%%%%%%%%%%%%%%%%%%%%
\subsection{Coordinate Charts}
\label{subsec:coordinate-charts-manifolds}
%%%%%%%%%%%%%%%%%%%%%%%%%%%%%%%%%%%%%%%%%%%%%%%%%%%%%%%%%%%%

\begin{definitionbox}{Coordinate Chart}
A \emph{coordinate chart} on an \(n\)-manifold \(M\) is a pair

\[
\boxed{
(U,\varphi)
}
\]

where \(U\subseteq M\) is open and

\[
\varphi:U\to\varphi(U)\subseteq\R^n
\]

is a homeomorphism onto an open subset of \(\R^n\).
\end{definitionbox}

The map

\[
\varphi
\]

assigns coordinates to points of \(U\).

%%%%%%%%%%%%%%%%%%%%%%%%%%%%%%%%%%%%%%%%%%%%%%%%%%%%%%%%%%%%
\subsection{Coordinate Functions}
\label{subsec:coordinate-functions}
%%%%%%%%%%%%%%%%%%%%%%%%%%%%%%%%%%%%%%%%%%%%%%%%%%%%%%%%%%%%

A chart may be written

\[
\varphi(p)
=
\left(
x^1(p),\ldots,x^n(p)
\right).
\]

The functions

\[
\boxed{
x^1,\ldots,x^n
}
\]

are called coordinate functions.

The superscripts are labels rather than exponents.

Thus

\[
x^2
\]

in this context means the second coordinate function, not necessarily the
square of \(x\).

%%%%%%%%%%%%%%%%%%%%%%%%%%%%%%%%%%%%%%%%%%%%%%%%%%%%%%%%%%%%
\subsection{Atlases}
\label{subsec:atlases-manifolds}
%%%%%%%%%%%%%%%%%%%%%%%%%%%%%%%%%%%%%%%%%%%%%%%%%%%%%%%%%%%%

\begin{definitionbox}{Atlas}
An \emph{atlas} on a manifold \(M\) is a collection of coordinate charts

\[
\boxed{
\{(U_\alpha,\varphi_\alpha)\}_{\alpha\in A}
}
\]

whose domains cover \(M\):

\[
\boxed{
M=\bigcup_{\alpha\in A}U_\alpha.
}
\]
\end{definitionbox}

The lowercase Greek letter

\[
\alpha
\]

is read ``alpha.''

It is commonly used as an index labeling charts.

%%%%%%%%%%%%%%%%%%%%%%%%%%%%%%%%%%%%%%%%%%%%%%%%%%%%%%%%%%%%
\subsection{Transition Maps}
\label{subsec:transition-maps-manifolds}
%%%%%%%%%%%%%%%%%%%%%%%%%%%%%%%%%%%%%%%%%%%%%%%%%%%%%%%%%%%%

Suppose two charts

\[
(U,\varphi)
\qquad\text{and}\qquad
(V,\psi)
\]

overlap.

On

\[
U\cap V,
\]

the corresponding coordinates are related by

\[
\boxed{
\psi\circ\varphi^{-1}.
}
\]

This map is called a \emph{transition map} or \emph{change of coordinates}.

It maps between open subsets of Euclidean space:

\[
\varphi(U\cap V)
\longrightarrow
\psi(U\cap V).
\]

%%%%%%%%%%%%%%%%%%%%%%%%%%%%%%%%%%%%%%%%%%%%%%%%%%%%%%%%%%%%
\subsection{Smooth Compatibility}
\label{subsec:smooth-compatibility-manifolds}
%%%%%%%%%%%%%%%%%%%%%%%%%%%%%%%%%%%%%%%%%%%%%%%%%%%%%%%%%%%%

Two charts are smoothly compatible when their transition maps are smooth
wherever the charts overlap.

That is,

\[
\boxed{
\psi\circ\varphi^{-1}
}
\]

and its inverse are smooth.

Since the transition map is already a homeomorphism, smoothness in both
directions means it is a diffeomorphism between open subsets of Euclidean
space.

%%%%%%%%%%%%%%%%%%%%%%%%%%%%%%%%%%%%%%%%%%%%%%%%%%%%%%%%%%%%
\subsection{Smooth Atlases}
\label{subsec:smooth-atlases}
%%%%%%%%%%%%%%%%%%%%%%%%%%%%%%%%%%%%%%%%%%%%%%%%%%%%%%%%%%%%

\begin{definitionbox}{Smooth Atlas}
A \emph{smooth atlas} is an atlas in which every pair of overlapping charts
is smoothly compatible.
\end{definitionbox}

A topological manifold equipped with a smooth atlas becomes a smooth
manifold.

%%%%%%%%%%%%%%%%%%%%%%%%%%%%%%%%%%%%%%%%%%%%%%%%%%%%%%%%%%%%
\subsection{Smooth Structures}
\label{subsec:manifold-smooth-structures}
%%%%%%%%%%%%%%%%%%%%%%%%%%%%%%%%%%%%%%%%%%%%%%%%%%%%%%%%%%%%

Two smooth atlases are considered equivalent when their union is again a
smooth atlas.

An equivalence class of compatible smooth atlases determines a
\emph{smooth structure} on \(M\).

A smooth structure can also be represented by a maximal smooth atlas.

%%%%%%%%%%%%%%%%%%%%%%%%%%%%%%%%%%%%%%%%%%%%%%%%%%%%%%%%%%%%
\subsection{Maximal Atlases}
\label{subsec:maximal-atlases}
%%%%%%%%%%%%%%%%%%%%%%%%%%%%%%%%%%%%%%%%%%%%%%%%%%%%%%%%%%%%

A maximal smooth atlas contains every chart smoothly compatible with the
given smooth structure.

One does not normally list all charts explicitly.

Instead, a convenient smaller atlas is given, and the corresponding maximal
atlas is understood.

%%%%%%%%%%%%%%%%%%%%%%%%%%%%%%%%%%%%%%%%%%%%%%%%%%%%%%%%%%%%
\subsection{Examples of Manifolds}
\label{subsec:examples-manifolds}
%%%%%%%%%%%%%%%%%%%%%%%%%%%%%%%%%%%%%%%%%%%%%%%%%%%%%%%%%%%%

Standard examples include

\[
\boxed{
\R^n,
}
\]

\[
\boxed{
S^n,
}
\]

\[
\boxed{
T^n,
}
\]

\[
\boxed{
\mathbb{RP}^n,
}
\]

\[
\boxed{
\mathbb{CP}^n,
}
\]

and matrix groups such as

\[
\boxed{
\mathrm{GL}(n,\R).
}
\]

Open subsets of manifolds are themselves manifolds of the same dimension.

%%%%%%%%%%%%%%%%%%%%%%%%%%%%%%%%%%%%%%%%%%%%%%%%%%%%%%%%%%%%
\subsection{Euclidean Space}
\label{subsec:euclidean-space-manifold}
%%%%%%%%%%%%%%%%%%%%%%%%%%%%%%%%%%%%%%%%%%%%%%%%%%%%%%%%%%%%

The simplest \(n\)-manifold is

\[
\R^n.
\]

A global coordinate chart is given by

\[
\boxed{
\operatorname{id}_{\R^n}:\R^n\to\R^n.
}
\]

Thus Euclidean space is both locally and globally Euclidean.

%%%%%%%%%%%%%%%%%%%%%%%%%%%%%%%%%%%%%%%%%%%%%%%%%%%%%%%%%%%%
\subsection{The Circle as a Manifold}
\label{subsec:circle-manifold}
%%%%%%%%%%%%%%%%%%%%%%%%%%%%%%%%%%%%%%%%%%%%%%%%%%%%%%%%%%%%

The circle

\[
S^1
=
\{(x,y)\in\R^2:x^2+y^2=1\}
\]

is not homeomorphic to \(\R\).

However, sufficiently small arcs are homeomorphic to open intervals.

Therefore

\[
\boxed{
S^1
\text{ is a }1\text{-manifold}.
}
\]

%%%%%%%%%%%%%%%%%%%%%%%%%%%%%%%%%%%%%%%%%%%%%%%%%%%%%%%%%%%%
\subsection{The Sphere as a Manifold}
\label{subsec:sphere-manifold}
%%%%%%%%%%%%%%%%%%%%%%%%%%%%%%%%%%%%%%%%%%%%%%%%%%%%%%%%%%%%

The \(n\)-sphere is

\[
\boxed{
S^n
=
\left\{
x\in\R^{n+1}:
\|x\|=1
\right\}.
}
\]

Each point has a neighborhood homeomorphic to an open subset of \(\R^n\).

Therefore

\[
\boxed{
S^n
\text{ is an }n\text{-manifold}.
}
\]

%%%%%%%%%%%%%%%%%%%%%%%%%%%%%%%%%%%%%%%%%%%%%%%%%%%%%%%%%%%%
\subsection{Stereographic Charts}
\label{subsec:stereographic-charts-manifolds}
%%%%%%%%%%%%%%%%%%%%%%%%%%%%%%%%%%%%%%%%%%%%%%%%%%%%%%%%%%%%

The sphere can be covered using stereographic projection.

Let

\[
N=(0,\ldots,0,1)
\]

and

\[
S=(0,\ldots,0,-1)
\]

be the north and south poles.

Then

\[
S^n\setminus\{N\}
\]

and

\[
S^n\setminus\{S\}
\]

are each diffeomorphic to

\[
\R^n.
\]

Thus two charts suffice to cover the sphere.

%%%%%%%%%%%%%%%%%%%%%%%%%%%%%%%%%%%%%%%%%%%%%%%%%%%%%%%%%%%%
\subsection{The Torus}
\label{subsec:torus-manifold}
%%%%%%%%%%%%%%%%%%%%%%%%%%%%%%%%%%%%%%%%%%%%%%%%%%%%%%%%%%%%

The \(n\)-torus is

\[
\boxed{
T^n
=
\underbrace{
S^1\times\cdots\times S^1
}_{n\text{ factors}}.
}
\]

Equivalently,

\[
\boxed{
T^n
\cong
\R^n/\Z^n.
}
\]

For \(n=2\),

\[
T^2=S^1\times S^1.
\]

The torus illustrates both product and quotient constructions of manifolds.

%%%%%%%%%%%%%%%%%%%%%%%%%%%%%%%%%%%%%%%%%%%%%%%%%%%%%%%%%%%%
\subsection{Product Manifolds}
\label{subsec:product-manifolds}
%%%%%%%%%%%%%%%%%%%%%%%%%%%%%%%%%%%%%%%%%%%%%%%%%%%%%%%%%%%%

If \(M\) is an \(m\)-manifold and \(N\) is an \(n\)-manifold, then

\[
\boxed{
M\times N
}
\]

is an \((m+n)\)-manifold.

Therefore,

\[
\boxed{
\dim(M\times N)
=
\dim M+\dim N.
}
\]

For example,

\[
S^1\times S^1
\]

has dimension

\[
1+1=2.
\]

%%%%%%%%%%%%%%%%%%%%%%%%%%%%%%%%%%%%%%%%%%%%%%%%%%%%%%%%%%%%
\subsection{Product Charts}
\label{subsec:product-charts}
%%%%%%%%%%%%%%%%%%%%%%%%%%%%%%%%%%%%%%%%%%%%%%%%%%%%%%%%%%%%

Suppose

\[
(U,\varphi)
\]

is a chart on \(M\) and

\[
(V,\psi)
\]

is a chart on \(N\).

Then

\[
\boxed{
(U\times V,\varphi\times\psi)
}
\]

is a chart on

\[
M\times N.
\]

Locally,

\[
\R^m\times\R^n
\cong
\R^{m+n}.
\]

%%%%%%%%%%%%%%%%%%%%%%%%%%%%%%%%%%%%%%%%%%%%%%%%%%%%%%%%%%%%
\subsection{Quotient Manifolds}
\label{subsec:quotient-manifolds}
%%%%%%%%%%%%%%%%%%%%%%%%%%%%%%%%%%%%%%%%%%%%%%%%%%%%%%%%%%%%

Many important manifolds arise by identifying points under a group action.

For example,

\[
T^n
=
\R^n/\Z^n.
\]

The integer lattice acts on \(\R^n\) by translations:

\[
\boxed{
x\mapsto x+k,
\qquad
k\in\Z^n.
}
\]

The quotient identifies points differing by lattice vectors.

%%%%%%%%%%%%%%%%%%%%%%%%%%%%%%%%%%%%%%%%%%%%%%%%%%%%%%%%%%%%
\subsection{Group Actions and Quotients}
\label{subsec:group-actions-quotient-manifolds}
%%%%%%%%%%%%%%%%%%%%%%%%%%%%%%%%%%%%%%%%%%%%%%%%%%%%%%%%%%%%

Let a group \(G\) act smoothly on a manifold \(M\).

Under suitable hypotheses, especially when the action is free and proper,
the quotient

\[
\boxed{
M/G
}
\]

inherits a smooth manifold structure.

Moreover,

\[
\boxed{
\dim(M/G)
=
\dim M-\dim G
}
\]

when \(G\) is a Lie group acting freely and properly.

%%%%%%%%%%%%%%%%%%%%%%%%%%%%%%%%%%%%%%%%%%%%%%%%%%%%%%%%%%%%
\subsection{Projective Spaces}
\label{subsec:projective-spaces-manifolds}
%%%%%%%%%%%%%%%%%%%%%%%%%%%%%%%%%%%%%%%%%%%%%%%%%%%%%%%%%%%%

Real projective space is

\[
\boxed{
\mathbb{RP}^n
=
S^n/\{x\sim -x\}.
}
\]

Thus antipodal points are identified.

Complex projective space is

\[
\boxed{
\mathbb{CP}^n
=
(\mathbb{C}^{n+1}\setminus\{0\})/\sim,
}
\]

where

\[
z\sim\lambda z
\]

for every

\[
\lambda\in\mathbb{C}\setminus\{0\}.
\]

Its real dimension is

\[
\boxed{
\dim_{\R}\mathbb{CP}^n=2n.
}
\]

%%%%%%%%%%%%%%%%%%%%%%%%%%%%%%%%%%%%%%%%%%%%%%%%%%%%%%%%%%%%
\subsection{Manifolds with Boundary}
\label{subsec:manifolds-with-boundary}
%%%%%%%%%%%%%%%%%%%%%%%%%%%%%%%%%%%%%%%%%%%%%%%%%%%%%%%%%%%%

Not every manifold is locally modeled exclusively on \(\R^n\).

A manifold with boundary may also locally resemble the closed half-space

\[
\boxed{
\mathbb{H}^n
=
\{
(x^1,\ldots,x^n)\in\R^n:
x^n\geq0
\}.
}
\]

Points modeled on

\[
x^n=0
\]

form the boundary.

%%%%%%%%%%%%%%%%%%%%%%%%%%%%%%%%%%%%%%%%%%%%%%%%%%%%%%%%%%%%
\subsection{Interior and Boundary}
\label{subsec:manifold-interior-boundary}
%%%%%%%%%%%%%%%%%%%%%%%%%%%%%%%%%%%%%%%%%%%%%%%%%%%%%%%%%%%%

For a manifold with boundary \(M\), one writes

\[
\boxed{
\operatorname{Int}(M)
}
\]

for the manifold interior and

\[
\boxed{
\partial M
}
\]

for its boundary.

The boundary of an \(n\)-manifold is an \((n-1)\)-manifold:

\[
\boxed{
\dim(\partial M)=n-1.
}
\]

%%%%%%%%%%%%%%%%%%%%%%%%%%%%%%%%%%%%%%%%%%%%%%%%%%%%%%%%%%%%
\subsection{Examples with Boundary}
\label{subsec:examples-manifolds-boundary}
%%%%%%%%%%%%%%%%%%%%%%%%%%%%%%%%%%%%%%%%%%%%%%%%%%%%%%%%%%%%

The interval

\[
[0,1]
\]

is a \(1\)-manifold with boundary

\[
\boxed{
\partial[0,1]=\{0,1\}.
}
\]

The disk

\[
D^2
=
\{(x,y):x^2+y^2\leq1\}
\]

is a \(2\)-manifold with boundary

\[
\boxed{
\partial D^2=S^1.
}
\]

More generally,

\[
\boxed{
\partial D^n=S^{n-1}.
}
\]

%%%%%%%%%%%%%%%%%%%%%%%%%%%%%%%%%%%%%%%%%%%%%%%%%%%%%%%%%%%%
\subsection{Boundary of a Boundary}
\label{subsec:boundary-boundary-manifolds}
%%%%%%%%%%%%%%%%%%%%%%%%%%%%%%%%%%%%%%%%%%%%%%%%%%%%%%%%%%%%

For an ordinary manifold with boundary,

\[
\boxed{
\partial(\partial M)=\varnothing.
}
\]

This geometric fact parallels the algebraic identity

\[
\boxed{
\partial^2=0
}
\]

from chain complexes and homology.

%%%%%%%%%%%%%%%%%%%%%%%%%%%%%%%%%%%%%%%%%%%%%%%%%%%%%%%%%%%%
\subsection{Submanifolds}
\label{subsec:submanifolds}
%%%%%%%%%%%%%%%%%%%%%%%%%%%%%%%%%%%%%%%%%%%%%%%%%%%%%%%%%%%%

A submanifold is a subset of a manifold that itself carries a compatible
manifold structure.

For example,

\[
\boxed{
S^n\subseteq\R^{n+1}
}
\]

is an embedded \(n\)-dimensional submanifold.

Likewise,

\[
\boxed{
T^2\subseteq\R^3
}
\]

can be realized as an embedded surface.

%%%%%%%%%%%%%%%%%%%%%%%%%%%%%%%%%%%%%%%%%%%%%%%%%%%%%%%%%%%%
\subsection{Regular Level Sets as Submanifolds}
\label{subsec:regular-level-set-submanifolds}
%%%%%%%%%%%%%%%%%%%%%%%%%%%%%%%%%%%%%%%%%%%%%%%%%%%%%%%%%%%%

Let

\[
F:M^m\to N^n
\]

be smooth.

If

\[
y\in N
\]

is a regular value, then

\[
\boxed{
F^{-1}(y)
}
\]

is a submanifold of dimension

\[
\boxed{
m-n.
}
\]

This is the Regular Value Theorem developed in Differential Topology.

%%%%%%%%%%%%%%%%%%%%%%%%%%%%%%%%%%%%%%%%%%%%%%%%%%%%%%%%%%%%
\subsection{Worked Example: The Sphere as a Manifold}
\label{subsec:worked-sphere-manifold}
%%%%%%%%%%%%%%%%%%%%%%%%%%%%%%%%%%%%%%%%%%%%%%%%%%%%%%%%%%%%

\begin{workedexamplebox}{Dimension of the Sphere}

Define

\[
F:\R^{n+1}\to\R
\]

by

\[
F(x)=\|x\|^2.
\]

Then

\[
S^n=F^{-1}(1).
\]

Since

\[
dF_x(v)=2x\cdot v,
\]

the differential is nonzero whenever

\[
x\in S^n.
\]

Thus \(1\) is a regular value.

Therefore,

\[
\dim S^n
=
(n+1)-1.
\]

\begin{answerbox}
\[
\boxed{
\dim S^n=n.
}
\]
\end{answerbox}
\end{workedexamplebox}

%%%%%%%%%%%%%%%%%%%%%%%%%%%%%%%%%%%%%%%%%%%%%%%%%%%%%%%%%%%%
\subsection{Smooth Maps Between Manifolds}
\label{subsec:smooth-maps-manifolds}
%%%%%%%%%%%%%%%%%%%%%%%%%%%%%%%%%%%%%%%%%%%%%%%%%%%%%%%%%%%%

A map

\[
F:M\to N
\]

between smooth manifolds is smooth when its coordinate representation is
smooth.

Given charts

\[
(U,\varphi)
\]

and

\[
(V,\psi),
\]

the coordinate representation is

\[
\boxed{
\psi\circ F\circ\varphi^{-1}.
}
\]

Smoothness is therefore checked using ordinary multivariable calculus in
local coordinates.

%%%%%%%%%%%%%%%%%%%%%%%%%%%%%%%%%%%%%%%%%%%%%%%%%%%%%%%%%%%%
\subsection{Diffeomorphisms}
\label{subsec:manifold-diffeomorphisms}
%%%%%%%%%%%%%%%%%%%%%%%%%%%%%%%%%%%%%%%%%%%%%%%%%%%%%%%%%%%%

A diffeomorphism is a smooth bijection

\[
F:M\to N
\]

whose inverse is also smooth.

If such a map exists, \(M\) and \(N\) have the same smooth-manifold
structure up to equivalence.

We write

\[
\boxed{
M\cong_{\mathrm{diff}}N.
}
\]

%%%%%%%%%%%%%%%%%%%%%%%%%%%%%%%%%%%%%%%%%%%%%%%%%%%%%%%%%%%%
\subsection{Differential of a Map}
\label{subsec:manifold-differential-map}
%%%%%%%%%%%%%%%%%%%%%%%%%%%%%%%%%%%%%%%%%%%%%%%%%%%%%%%%%%%%

A smooth map

\[
F:M\to N
\]

induces at each point \(p\in M\) a linear map

\[
\boxed{
dF_p:T_pM\to T_{F(p)}N.
}
\]

The differential generalizes the derivative and Jacobian from Euclidean
calculus.

%%%%%%%%%%%%%%%%%%%%%%%%%%%%%%%%%%%%%%%%%%%%%%%%%%%%%%%%%%%%
\subsection{Tangent Spaces}
\label{subsec:tangent-spaces-manifolds}
%%%%%%%%%%%%%%%%%%%%%%%%%%%%%%%%%%%%%%%%%%%%%%%%%%%%%%%%%%%%

The tangent space at \(p\in M\) is denoted

\[
\boxed{
T_pM.
}
\]

It is an \(n\)-dimensional vector space when

\[
\dim M=n.
\]

Tangent vectors can be defined in several equivalent ways, including:

\begin{itemize}
    \item velocities of curves;
    \item derivations acting on smooth functions;
    \item coordinate-vector equivalence classes.
\end{itemize}

%%%%%%%%%%%%%%%%%%%%%%%%%%%%%%%%%%%%%%%%%%%%%%%%%%%%%%%%%%%%
\subsection{Tangent Vectors from Curves}
\label{subsec:tangent-vectors-curves}
%%%%%%%%%%%%%%%%%%%%%%%%%%%%%%%%%%%%%%%%%%%%%%%%%%%%%%%%%%%%

Let

\[
\gamma:(-\varepsilon,\varepsilon)\to M
\]

be a smooth curve with

\[
\gamma(0)=p.
\]

The velocity

\[
\boxed{
\gamma'(0)
}
\]

determines a tangent vector in

\[
T_pM.
\]

Two curves determine the same tangent vector when they have the same
first-order behavior in local coordinates.

%%%%%%%%%%%%%%%%%%%%%%%%%%%%%%%%%%%%%%%%%%%%%%%%%%%%%%%%%%%%
\subsection{Coordinate Basis of a Tangent Space}
\label{subsec:coordinate-basis-tangent}
%%%%%%%%%%%%%%%%%%%%%%%%%%%%%%%%%%%%%%%%%%%%%%%%%%%%%%%%%%%%

Given local coordinates

\[
(x^1,\ldots,x^n),
\]

the tangent space has coordinate basis

\[
\boxed{
\left\{
\frac{\partial}{\partial x^1}\bigg|_p,
\ldots,
\frac{\partial}{\partial x^n}\bigg|_p
\right\}.
}
\]

A tangent vector can therefore be written

\[
\boxed{
v
=
v^i
\frac{\partial}{\partial x^i}\bigg|_p,
}
\]

using Einstein summation notation.

%%%%%%%%%%%%%%%%%%%%%%%%%%%%%%%%%%%%%%%%%%%%%%%%%%%%%%%%%%%%
\subsection{Tangent Bundle}
\label{subsec:tangent-bundle-manifolds}
%%%%%%%%%%%%%%%%%%%%%%%%%%%%%%%%%%%%%%%%%%%%%%%%%%%%%%%%%%%%

The tangent spaces can be assembled into a single space.

\begin{definitionbox}{Tangent Bundle}
The \emph{tangent bundle} of \(M\) is

\[
\boxed{
TM
=
\bigsqcup_{p\in M}T_pM.
}
\]
\end{definitionbox}

There is a natural projection

\[
\boxed{
\pi:TM\to M
}
\]

defined by

\[
\pi(v)=p
\]

when

\[
v\in T_pM.
\]

If

\[
\dim M=n,
\]

then

\[
\boxed{
\dim TM=2n.
}
\]

%%%%%%%%%%%%%%%%%%%%%%%%%%%%%%%%%%%%%%%%%%%%%%%%%%%%%%%%%%%%
\subsection{Cotangent Spaces}
\label{subsec:cotangent-spaces-manifolds}
%%%%%%%%%%%%%%%%%%%%%%%%%%%%%%%%%%%%%%%%%%%%%%%%%%%%%%%%%%%%

The dual vector space of \(T_pM\) is the cotangent space

\[
\boxed{
T_p^*M.
}
\]

Its elements are linear maps

\[
\boxed{
\omega:T_pM\to\R.
}
\]

In local coordinates, the natural dual basis is

\[
\boxed{
dx^1,\ldots,dx^n.
}
\]

%%%%%%%%%%%%%%%%%%%%%%%%%%%%%%%%%%%%%%%%%%%%%%%%%%%%%%%%%%%%
\subsection{Cotangent Bundle}
\label{subsec:cotangent-bundle-manifolds}
%%%%%%%%%%%%%%%%%%%%%%%%%%%%%%%%%%%%%%%%%%%%%%%%%%%%%%%%%%%%

The cotangent bundle is

\[
\boxed{
T^*M
=
\bigsqcup_{p\in M}T_p^*M.
}
\]

Like the tangent bundle,

\[
\boxed{
\dim T^*M=2n
}
\]

when

\[
\dim M=n.
\]

Cotangent bundles are fundamental in differential geometry, symplectic
geometry, and Hamiltonian mechanics.

%%%%%%%%%%%%%%%%%%%%%%%%%%%%%%%%%%%%%%%%%%%%%%%%%%%%%%%%%%%%
\subsection{Vector Fields}
\label{subsec:vector-fields-manifolds}
%%%%%%%%%%%%%%%%%%%%%%%%%%%%%%%%%%%%%%%%%%%%%%%%%%%%%%%%%%%%

A vector field is a smooth assignment

\[
\boxed{
p\mapsto X_p\in T_pM.
}
\]

Equivalently, it is a smooth section

\[
X:M\to TM
\]

satisfying

\[
\boxed{
\pi\circ X=\operatorname{id}_M.
}
\]

In local coordinates,

\[
\boxed{
X
=
X^i
\frac{\partial}{\partial x^i}.
}
\]

%%%%%%%%%%%%%%%%%%%%%%%%%%%%%%%%%%%%%%%%%%%%%%%%%%%%%%%%%%%%
\subsection{Differential Forms}
\label{subsec:differential-forms-manifolds}
%%%%%%%%%%%%%%%%%%%%%%%%%%%%%%%%%%%%%%%%%%%%%%%%%%%%%%%%%%%%

A differential \(1\)-form assigns a cotangent vector to each point.

More generally, a differential \(k\)-form assigns an alternating
\(k\)-linear map on tangent vectors.

The space of smooth \(k\)-forms is commonly denoted

\[
\boxed{
\Omega^k(M).
}
\]

The uppercase Greek letter

\[
\Omega
\]

is read ``Omega.''

Differential forms provide the natural language for integration on
manifolds.

%%%%%%%%%%%%%%%%%%%%%%%%%%%%%%%%%%%%%%%%%%%%%%%%%%%%%%%%%%%%
\subsection{Immersions}
\label{subsec:manifold-immersions}
%%%%%%%%%%%%%%%%%%%%%%%%%%%%%%%%%%%%%%%%%%%%%%%%%%%%%%%%%%%%

A smooth map

\[
F:M\to N
\]

is an immersion when

\[
\boxed{
dF_p
}
\]

is injective for every \(p\in M\).

An immersion preserves the local dimension of tangent directions but may
have global self-intersections.

%%%%%%%%%%%%%%%%%%%%%%%%%%%%%%%%%%%%%%%%%%%%%%%%%%%%%%%%%%%%
\subsection{Submersions}
\label{subsec:manifold-submersions}
%%%%%%%%%%%%%%%%%%%%%%%%%%%%%%%%%%%%%%%%%%%%%%%%%%%%%%%%%%%%

A smooth map

\[
F:M\to N
\]

is a submersion when

\[
\boxed{
dF_p
}
\]

is surjective for every \(p\in M\).

Submersions locally behave like coordinate projections.

%%%%%%%%%%%%%%%%%%%%%%%%%%%%%%%%%%%%%%%%%%%%%%%%%%%%%%%%%%%%
\subsection{Embeddings}
\label{subsec:manifold-embeddings}
%%%%%%%%%%%%%%%%%%%%%%%%%%%%%%%%%%%%%%%%%%%%%%%%%%%%%%%%%%%%

A smooth embedding is an immersion

\[
F:M\to N
\]

that is also a homeomorphism onto its image.

Thus

\[
\boxed{
\text{embedding}
=
\text{local smooth inclusion}
+
\text{correct global topology}.
}
\]

%%%%%%%%%%%%%%%%%%%%%%%%%%%%%%%%%%%%%%%%%%%%%%%%%%%%%%%%%%%%
\subsection{Whitney Embedding Theorem}
\label{subsec:manifold-whitney-embedding}
%%%%%%%%%%%%%%%%%%%%%%%%%%%%%%%%%%%%%%%%%%%%%%%%%%%%%%%%%%%%

\begin{theorem}[Whitney Embedding Theorem]
Every smooth \(n\)-manifold satisfying the standard manifold assumptions can
be smoothly embedded in

\[
\boxed{
\R^{2n}.
}
\]
\end{theorem}

Therefore every abstract smooth manifold can be realized as a concrete
submanifold of some Euclidean space.

%%%%%%%%%%%%%%%%%%%%%%%%%%%%%%%%%%%%%%%%%%%%%%%%%%%%%%%%%%%%
\subsection{Orientability}
\label{subsec:manifold-orientability}
%%%%%%%%%%%%%%%%%%%%%%%%%%%%%%%%%%%%%%%%%%%%%%%%%%%%%%%%%%%%

An orientation provides a consistent choice of orientation for tangent
spaces.

In coordinates, orientation-preserving transition maps satisfy

\[
\boxed{
\det
\left(
\frac{\partial y^i}{\partial x^j}
\right)
>0.
}
\]

A manifold that admits such a consistent choice is called orientable.

%%%%%%%%%%%%%%%%%%%%%%%%%%%%%%%%%%%%%%%%%%%%%%%%%%%%%%%%%%%%
\subsection{Orientation Forms}
\label{subsec:orientation-forms-manifolds}
%%%%%%%%%%%%%%%%%%%%%%%%%%%%%%%%%%%%%%%%%%%%%%%%%%%%%%%%%%%%

For a smooth \(n\)-manifold, orientability is equivalent to the existence of
a nowhere-vanishing smooth \(n\)-form.

Thus there exists

\[
\boxed{
\omega\in\Omega^n(M)
}
\]

such that

\[
\boxed{
\omega_p\neq0
}
\]

for every \(p\in M\).

Such a form is called a volume form or orientation form.

%%%%%%%%%%%%%%%%%%%%%%%%%%%%%%%%%%%%%%%%%%%%%%%%%%%%%%%%%%%%
\subsection{Examples of Orientability}
\label{subsec:examples-orientability-manifolds}
%%%%%%%%%%%%%%%%%%%%%%%%%%%%%%%%%%%%%%%%%%%%%%%%%%%%%%%%%%%%

The following manifolds are orientable:

\[
\boxed{
\R^n,
\quad
S^n,
\quad
T^n.
}
\]

The M\"obius strip is nonorientable.

The real projective plane

\[
\boxed{
\mathbb{RP}^2
}
\]

is also nonorientable.

More generally,

\[
\boxed{
\mathbb{RP}^n
\text{ is orientable exactly when }n\text{ is odd}.
}
\]

%%%%%%%%%%%%%%%%%%%%%%%%%%%%%%%%%%%%%%%%%%%%%%%%%%%%%%%%%%%%
\subsection{Partitions of Unity}
\label{subsec:partitions-of-unity-manifolds}
%%%%%%%%%%%%%%%%%%%%%%%%%%%%%%%%%%%%%%%%%%%%%%%%%%%%%%%%%%%%

A major challenge on manifolds is turning local constructions into global
ones.

Partitions of unity solve this problem.

\begin{definitionbox}{Partition of Unity}
A smooth partition of unity subordinate to an open cover
\(\{U_\alpha\}\) is a collection of smooth functions

\[
\boxed{
\rho_\alpha:M\to[0,1]
}
\]

such that:

\begin{enumerate}
    \item the family of supports is locally finite;
    \item \(\operatorname{supp}\rho_\alpha\subseteq U_\alpha\);
    \item
    \[
    \boxed{
    \sum_\alpha\rho_\alpha(p)=1
    }
    \]
    for every \(p\in M\).
\end{enumerate}
\end{definitionbox}

The lowercase Greek letter

\[
\rho
\]

is read ``rho.''

%%%%%%%%%%%%%%%%%%%%%%%%%%%%%%%%%%%%%%%%%%%%%%%%%%%%%%%%%%%%
\subsection{Why Partitions of Unity Matter}
\label{subsec:why-partitions-unity}
%%%%%%%%%%%%%%%%%%%%%%%%%%%%%%%%%%%%%%%%%%%%%%%%%%%%%%%%%%%%

Suppose local objects

\[
A_\alpha
\]

are defined on charts \(U_\alpha\).

A partition of unity allows one to combine them schematically as

\[
\boxed{
A
=
\sum_\alpha\rho_\alpha A_\alpha.
}
\]

This technique is used to construct global

\[
\boxed{
\text{metrics},
\quad
\text{connections},
\quad
\text{functions},
\quad
\text{differential forms}.
}
\]

%%%%%%%%%%%%%%%%%%%%%%%%%%%%%%%%%%%%%%%%%%%%%%%%%%%%%%%%%%%%
\subsection{Existence of Partitions of Unity}
\label{subsec:existence-partitions-unity}
%%%%%%%%%%%%%%%%%%%%%%%%%%%%%%%%%%%%%%%%%%%%%%%%%%%%%%%%%%%%

Smooth manifolds under the standard Hausdorff and second-countability
assumptions are paracompact.

Consequently, suitable open covers admit smooth partitions of unity
subordinate to refinements of those covers.

This theorem is one of the most important local-to-global tools in manifold
theory.

%%%%%%%%%%%%%%%%%%%%%%%%%%%%%%%%%%%%%%%%%%%%%%%%%%%%%%%%%%%%
\subsection{Riemannian Metrics}
\label{subsec:riemannian-metrics-manifolds}
%%%%%%%%%%%%%%%%%%%%%%%%%%%%%%%%%%%%%%%%%%%%%%%%%%%%%%%%%%%%

A smooth manifold initially has no intrinsic notion of length or angle.

A Riemannian metric adds such structure.

At each point \(p\), a Riemannian metric assigns an inner product

\[
\boxed{
g_p:T_pM\times T_pM\to\R
}
\]

that varies smoothly with \(p\).

%%%%%%%%%%%%%%%%%%%%%%%%%%%%%%%%%%%%%%%%%%%%%%%%%%%%%%%%%%%%
\subsection{Existence of Riemannian Metrics}
\label{subsec:existence-riemannian-metrics}
%%%%%%%%%%%%%%%%%%%%%%%%%%%%%%%%%%%%%%%%%%%%%%%%%%%%%%%%%%%%

Every smooth manifold under the standard assumptions admits a Riemannian
metric.

The construction uses local Euclidean inner products and a partition of
unity.

Thus

\[
\boxed{
\text{smooth manifold}
\Longrightarrow
\text{can be equipped with Riemannian geometry}.
}
\]

The metric is generally not unique.

%%%%%%%%%%%%%%%%%%%%%%%%%%%%%%%%%%%%%%%%%%%%%%%%%%%%%%%%%%%%
\subsection{Lie Groups}
\label{subsec:lie-groups-manifolds}
%%%%%%%%%%%%%%%%%%%%%%%%%%%%%%%%%%%%%%%%%%%%%%%%%%%%%%%%%%%%

A Lie group combines manifold structure with group structure.

\begin{definitionbox}{Lie Group}
A \emph{Lie group} is a smooth manifold \(G\) that is also a group such that
multiplication

\[
\boxed{
G\times G\to G,
\qquad
(g,h)\mapsto gh,
}
\]

and inversion

\[
\boxed{
G\to G,
\qquad
g\mapsto g^{-1},
}
\]

are smooth.
\end{definitionbox}

Examples include

\[
\boxed{
\R^n,
\quad
S^1,
\quad
\mathrm{GL}(n,\R),
\quad
\mathrm{SO}(n).
}
\]

%%%%%%%%%%%%%%%%%%%%%%%%%%%%%%%%%%%%%%%%%%%%%%%%%%%%%%%%%%%%
\subsection{General Linear Group as a Manifold}
\label{subsec:general-linear-manifold}
%%%%%%%%%%%%%%%%%%%%%%%%%%%%%%%%%%%%%%%%%%%%%%%%%%%%%%%%%%%%

The general linear group is

\[
\boxed{
\mathrm{GL}(n,\R)
=
\{
A\in M_n(\R):\det A\neq0
\}.
}
\]

Since the determinant is continuous,

\[
\det^{-1}(\R\setminus\{0\})
\]

is open.

Therefore

\[
\mathrm{GL}(n,\R)
\]

is an open subset of

\[
M_n(\R)\cong\R^{n^2}.
\]

Hence

\[
\boxed{
\dim\mathrm{GL}(n,\R)=n^2.
}
\]

%%%%%%%%%%%%%%%%%%%%%%%%%%%%%%%%%%%%%%%%%%%%%%%%%%%%%%%%%%%%
\subsection{Fiber Bundles}
\label{subsec:fiber-bundles-manifolds}
%%%%%%%%%%%%%%%%%%%%%%%%%%%%%%%%%%%%%%%%%%%%%%%%%%%%%%%%%%%%

Many manifolds naturally occur as families of spaces parametrized by
another manifold.

\begin{definitionbox}{Fiber Bundle}
A \emph{fiber bundle} consists of a space \(E\), a base space \(B\), a
projection

\[
\boxed{
\pi:E\to B,
}
\]

and a typical fiber \(F\), such that every point of \(B\) has a neighborhood
\(U\) for which

\[
\boxed{
\pi^{-1}(U)\cong U\times F
}
\]

in a manner compatible with the projection onto \(U\).
\end{definitionbox}

The space \(E\) is called the total space.

%%%%%%%%%%%%%%%%%%%%%%%%%%%%%%%%%%%%%%%%%%%%%%%%%%%%%%%%%%%%
\subsection{Local Triviality}
\label{subsec:local-triviality-bundles}
%%%%%%%%%%%%%%%%%%%%%%%%%%%%%%%%%%%%%%%%%%%%%%%%%%%%%%%%%%%%

The defining feature of a fiber bundle is local triviality.

Locally,

\[
\boxed{
E\sim B\times F.
}
\]

Globally, however, the bundle may be twisted.

This is another example of the recurring manifold principle:

\[
\boxed{
\text{locally simple}
\qquad
\text{globally nontrivial}.
}
\]

%%%%%%%%%%%%%%%%%%%%%%%%%%%%%%%%%%%%%%%%%%%%%%%%%%%%%%%%%%%%
\subsection{The M\"obius Bundle}
\label{subsec:mobius-bundle}
%%%%%%%%%%%%%%%%%%%%%%%%%%%%%%%%%%%%%%%%%%%%%%%%%%%%%%%%%%%%

The M\"obius strip can be viewed as a line bundle over

\[
S^1.
\]

Locally it looks like

\[
\boxed{
U\times\R.
}
\]

Globally, however, the fibers twist once around the circle.

It is therefore not globally equivalent to the trivial bundle

\[
S^1\times\R.
\]

%%%%%%%%%%%%%%%%%%%%%%%%%%%%%%%%%%%%%%%%%%%%%%%%%%%%%%%%%%%%
\subsection{Vector Bundles}
\label{subsec:vector-bundles-manifolds}
%%%%%%%%%%%%%%%%%%%%%%%%%%%%%%%%%%%%%%%%%%%%%%%%%%%%%%%%%%%%

A vector bundle is a fiber bundle whose fibers are vector spaces and whose
local trivializations preserve the linear structure.

A rank-\(k\) real vector bundle has typical fiber

\[
\boxed{
\R^k.
}
\]

The tangent bundle

\[
TM
\]

is a rank-\(n\) vector bundle over an \(n\)-manifold.

%%%%%%%%%%%%%%%%%%%%%%%%%%%%%%%%%%%%%%%%%%%%%%%%%%%%%%%%%%%%
\subsection{Sections}
\label{subsec:bundle-sections}
%%%%%%%%%%%%%%%%%%%%%%%%%%%%%%%%%%%%%%%%%%%%%%%%%%%%%%%%%%%%

A section of a bundle

\[
\pi:E\to M
\]

is a map

\[
\boxed{
s:M\to E
}
\]

such that

\[
\boxed{
\pi\circ s=\operatorname{id}_M.
}
\]

For the tangent bundle, sections are vector fields.

For the cotangent bundle, sections are differential \(1\)-forms.

%%%%%%%%%%%%%%%%%%%%%%%%%%%%%%%%%%%%%%%%%%%%%%%%%%%%%%%%%%%%
\subsection{Trivial Bundles}
\label{subsec:trivial-bundles}
%%%%%%%%%%%%%%%%%%%%%%%%%%%%%%%%%%%%%%%%%%%%%%%%%%%%%%%%%%%%

A bundle is trivial if it is globally equivalent to a product:

\[
\boxed{
E\cong B\times F.
}
\]

For a rank-\(k\) vector bundle,

\[
\boxed{
E\cong B\times\R^k.
}
\]

A vector bundle is trivial precisely when it admits a global frame of
linearly independent sections.

%%%%%%%%%%%%%%%%%%%%%%%%%%%%%%%%%%%%%%%%%%%%%%%%%%%%%%%%%%%%
\subsection{Parallelizable Manifolds}
\label{subsec:parallelizable-manifolds}
%%%%%%%%%%%%%%%%%%%%%%%%%%%%%%%%%%%%%%%%%%%%%%%%%%%%%%%%%%%%

\begin{definitionbox}{Parallelizable Manifold}
A smooth \(n\)-manifold \(M\) is \emph{parallelizable} if its tangent bundle
is trivial:

\[
\boxed{
TM\cong M\times\R^n.
}
\]
\end{definitionbox}

Examples include

\[
\boxed{
T^n
}
\]

and all Lie groups.

Among spheres, the parallelizable spheres are

\[
\boxed{
S^0,\quad S^1,\quad S^3,\quad S^7.
}
\]

%%%%%%%%%%%%%%%%%%%%%%%%%%%%%%%%%%%%%%%%%%%%%%%%%%%%%%%%%%%%
\subsection{Covering Spaces as Bundles}
\label{subsec:covering-spaces-bundles}
%%%%%%%%%%%%%%%%%%%%%%%%%%%%%%%%%%%%%%%%%%%%%%%%%%%%%%%%%%%%

A covering map

\[
p:E\to B
\]

can be viewed as a fiber bundle with discrete fiber.

Locally,

\[
\boxed{
p^{-1}(U)
}
\]

is a disjoint union of copies of \(U\).

For example,

\[
\boxed{
\R\to S^1,
\qquad
t\mapsto e^{2\pi i t},
}
\]

is a covering map with fiber

\[
\Z.
\]

%%%%%%%%%%%%%%%%%%%%%%%%%%%%%%%%%%%%%%%%%%%%%%%%%%%%%%%%%%%%
\subsection{Principal Bundles}
\label{subsec:principal-bundles}
%%%%%%%%%%%%%%%%%%%%%%%%%%%%%%%%%%%%%%%%%%%%%%%%%%%%%%%%%%%%

A principal bundle is a fiber bundle equipped with a compatible action of a
Lie group \(G\).

Its typical fiber is \(G\) itself.

Principal bundles provide the natural framework for

\[
\boxed{
\text{connections},
\quad
\text{gauge theory},
\quad
\text{characteristic classes}.
}
\]

%%%%%%%%%%%%%%%%%%%%%%%%%%%%%%%%%%%%%%%%%%%%%%%%%%%%%%%%%%%%
\subsection{Frame Bundle}
\label{subsec:frame-bundle}
%%%%%%%%%%%%%%%%%%%%%%%%%%%%%%%%%%%%%%%%%%%%%%%%%%%%%%%%%%%%

At each point \(p\in M\), consider all ordered bases of

\[
T_pM.
\]

These bases assemble into the frame bundle

\[
\boxed{
\operatorname{Fr}(M).
}
\]

For an \(n\)-manifold, this is naturally a principal

\[
\boxed{
\mathrm{GL}(n,\R)
}
\]

bundle.

Additional geometric structures can reduce the relevant structure group.

%%%%%%%%%%%%%%%%%%%%%%%%%%%%%%%%%%%%%%%%%%%%%%%%%%%%%%%%%%%%
\subsection{Characteristic Classes}
\label{subsec:characteristic-classes-manifolds}
%%%%%%%%%%%%%%%%%%%%%%%%%%%%%%%%%%%%%%%%%%%%%%%%%%%%%%%%%%%%

Characteristic classes associate cohomology classes to vector bundles.

They measure global twisting that cannot necessarily be detected in a
single coordinate chart.

Important examples include

\[
\boxed{
\text{Stiefel--Whitney classes},
}
\]

\[
\boxed{
\text{Chern classes},
}
\]

\[
\boxed{
\text{Pontryagin classes},
}
\]

and

\[
\boxed{
\text{Euler class}.
}
\]

%%%%%%%%%%%%%%%%%%%%%%%%%%%%%%%%%%%%%%%%%%%%%%%%%%%%%%%%%%%%
\subsection{Stiefel--Whitney Classes}
\label{subsec:stiefel-whitney-manifolds}
%%%%%%%%%%%%%%%%%%%%%%%%%%%%%%%%%%%%%%%%%%%%%%%%%%%%%%%%%%%%

For a real vector bundle \(E\), its Stiefel--Whitney classes are

\[
\boxed{
w_i(E)
\in
H^i(M;\Z/2\Z).
}
\]

The first class

\[
\boxed{
w_1(E)
}
\]

detects orientability in important cases.

For the tangent bundle,

\[
\boxed{
w_1(TM)=0
}
\]

if and only if \(M\) is orientable.

%%%%%%%%%%%%%%%%%%%%%%%%%%%%%%%%%%%%%%%%%%%%%%%%%%%%%%%%%%%%
\subsection{Chern Classes}
\label{subsec:chern-classes-manifolds}
%%%%%%%%%%%%%%%%%%%%%%%%%%%%%%%%%%%%%%%%%%%%%%%%%%%%%%%%%%%%

For a complex vector bundle \(E\), the Chern classes are

\[
\boxed{
c_i(E)
\in
H^{2i}(M;\Z).
}
\]

They play fundamental roles in

\[
\boxed{
\text{complex geometry},
\quad
\text{algebraic geometry},
\quad
\text{topology},
\quad
\text{mathematical physics}.
}
\]

%%%%%%%%%%%%%%%%%%%%%%%%%%%%%%%%%%%%%%%%%%%%%%%%%%%%%%%%%%%%
\subsection{Euler Class}
\label{subsec:euler-class-manifolds}
%%%%%%%%%%%%%%%%%%%%%%%%%%%%%%%%%%%%%%%%%%%%%%%%%%%%%%%%%%%%

For an oriented rank-\(n\) real vector bundle, the Euler class lies in

\[
\boxed{
e(E)\in H^n(M;\Z).
}
\]

For tangent bundles, the Euler class is closely connected to the Euler
characteristic.

For a closed oriented \(n\)-manifold,

\[
\boxed{
\langle e(TM),[M]\rangle
=
\chi(M).
}
\]

Here

\[
[M]
\]

is the fundamental class.

%%%%%%%%%%%%%%%%%%%%%%%%%%%%%%%%%%%%%%%%%%%%%%%%%%%%%%%%%%%%
\subsection{Compact Manifolds}
\label{subsec:compact-manifolds}
%%%%%%%%%%%%%%%%%%%%%%%%%%%%%%%%%%%%%%%%%%%%%%%%%%%%%%%%%%%%

A manifold may be compact or noncompact.

Examples of compact manifolds include

\[
\boxed{
S^n,
\quad
T^n,
\quad
\mathbb{RP}^n.
}
\]

Euclidean space

\[
\R^n
\]

is noncompact.

Compactness has powerful consequences because global arguments can often be
reduced to finite data.

%%%%%%%%%%%%%%%%%%%%%%%%%%%%%%%%%%%%%%%%%%%%%%%%%%%%%%%%%%%%
\subsection{Closed Manifolds}
\label{subsec:closed-manifolds}
%%%%%%%%%%%%%%%%%%%%%%%%%%%%%%%%%%%%%%%%%%%%%%%%%%%%%%%%%%%%

In manifold theory, a \emph{closed manifold} means a manifold that is

\[
\boxed{
\text{compact}
\quad\text{and}\quad
\text{without boundary}.
}
\]

This use of the word ``closed'' is stronger than merely being a closed
subset of some ambient space.

For example,

\[
\boxed{
S^n
}
\]

is a closed manifold.

%%%%%%%%%%%%%%%%%%%%%%%%%%%%%%%%%%%%%%%%%%%%%%%%%%%%%%%%%%%%
\subsection{Connected Manifolds}
\label{subsec:connected-manifolds}
%%%%%%%%%%%%%%%%%%%%%%%%%%%%%%%%%%%%%%%%%%%%%%%%%%%%%%%%%%%%

A connected manifold cannot be written as a union of two disjoint nonempty
open subsets.

For manifolds, connectedness has an especially useful consequence:

\[
\boxed{
\text{connected}
\Longrightarrow
\text{path connected}.
}
\]

Indeed, manifolds are locally path connected, so connected components and
path components coincide.

%%%%%%%%%%%%%%%%%%%%%%%%%%%%%%%%%%%%%%%%%%%%%%%%%%%%%%%%%%%%
\subsection{Components of a Manifold}
\label{subsec:components-manifold}
%%%%%%%%%%%%%%%%%%%%%%%%%%%%%%%%%%%%%%%%%%%%%%%%%%%%%%%%%%%%

Every connected component of a manifold is open.

Therefore a manifold is a disjoint union of connected manifolds:

\[
\boxed{
M
=
\bigsqcup_{\alpha}M_\alpha.
}
\]

If \(M\) is second countable, it has at most countably many connected
components.

%%%%%%%%%%%%%%%%%%%%%%%%%%%%%%%%%%%%%%%%%%%%%%%%%%%%%%%%%%%%
\subsection{Simply Connected Manifolds}
\label{subsec:simply-connected-manifolds}
%%%%%%%%%%%%%%%%%%%%%%%%%%%%%%%%%%%%%%%%%%%%%%%%%%%%%%%%%%%%

A connected manifold \(M\) is simply connected if

\[
\boxed{
\pi_1(M)=0.
}
\]

Examples include

\[
\boxed{
\R^n
}
\]

and, for \(n\geq2\),

\[
\boxed{
S^n.
}
\]

The torus is not simply connected:

\[
\boxed{
\pi_1(T^n)\cong\Z^n.
}
\]

%%%%%%%%%%%%%%%%%%%%%%%%%%%%%%%%%%%%%%%%%%%%%%%%%%%%%%%%%%%%
\subsection{Universal Covers of Manifolds}
\label{subsec:universal-covers-manifolds}
%%%%%%%%%%%%%%%%%%%%%%%%%%%%%%%%%%%%%%%%%%%%%%%%%%%%%%%%%%%%

A connected manifold is locally path connected and semilocally simply
connected.

Therefore it admits a universal covering space.

For example,

\[
\boxed{
\R^n\to T^n
}
\]

is the universal covering map of the \(n\)-torus.

The deck transformation group is

\[
\boxed{
\Z^n.
}
\]

%%%%%%%%%%%%%%%%%%%%%%%%%%%%%%%%%%%%%%%%%%%%%%%%%%%%%%%%%%%%
\subsection{Homology of Manifolds}
\label{subsec:homology-manifolds}
%%%%%%%%%%%%%%%%%%%%%%%%%%%%%%%%%%%%%%%%%%%%%%%%%%%%%%%%%%%%

Manifolds carry strong homological structure.

For a closed connected orientable \(n\)-manifold,

\[
\boxed{
H_n(M;\Z)\cong\Z.
}
\]

The generator is the fundamental class

\[
\boxed{
[M].
}
\]

For a closed connected nonorientable \(n\)-manifold,

\[
\boxed{
H_n(M;\Z)=0.
}
\]

%%%%%%%%%%%%%%%%%%%%%%%%%%%%%%%%%%%%%%%%%%%%%%%%%%%%%%%%%%%%
\subsection{Poincar\'e Duality}
\label{subsec:poincare-duality-manifolds}
%%%%%%%%%%%%%%%%%%%%%%%%%%%%%%%%%%%%%%%%%%%%%%%%%%%%%%%%%%%%

One of the deepest structural properties of closed oriented manifolds is
Poincar\'e duality.

\begin{theorem}[Poincar\'e Duality]
Let \(M\) be a closed connected oriented \(n\)-manifold.

Then, under the standard coefficient hypotheses,

\[
\boxed{
H^k(M)
\cong
H_{n-k}(M).
}
\]
\end{theorem}

More precisely, the isomorphism is induced by the fundamental class through
the cap product.

Poincar\'e duality reflects the fact that complementary-dimensional
topological features interact through intersection.

%%%%%%%%%%%%%%%%%%%%%%%%%%%%%%%%%%%%%%%%%%%%%%%%%%%%%%%%%%%%
\subsection{Manifolds and CW Complexes}
\label{subsec:manifolds-cw-complexes}
%%%%%%%%%%%%%%%%%%%%%%%%%%%%%%%%%%%%%%%%%%%%%%%%%%%%%%%%%%%%

Smooth manifolds have the homotopy type of CW complexes.

Compact smooth manifolds have the homotopy type of finite CW complexes.

This permits tools from algebraic topology to be applied systematically to
manifolds.

Morse theory provides one geometric mechanism for producing CW-type
decompositions.

%%%%%%%%%%%%%%%%%%%%%%%%%%%%%%%%%%%%%%%%%%%%%%%%%%%%%%%%%%%%
\subsection{Triangulations of Manifolds}
\label{subsec:triangulations-manifolds}
%%%%%%%%%%%%%%%%%%%%%%%%%%%%%%%%%%%%%%%%%%%%%%%%%%%%%%%%%%%%

Smooth manifolds admit compatible triangulations.

Thus a smooth manifold can be studied using simplicial methods.

For arbitrary topological manifolds, triangulability becomes substantially
more subtle and is dimension dependent.

Therefore one should not assume without qualification that every
topological manifold admits a triangulation.

%%%%%%%%%%%%%%%%%%%%%%%%%%%%%%%%%%%%%%%%%%%%%%%%%%%%%%%%%%%%
\subsection{Topological, PL, and Smooth Manifolds}
\label{subsec:topological-pl-smooth-manifolds}
%%%%%%%%%%%%%%%%%%%%%%%%%%%%%%%%%%%%%%%%%%%%%%%%%%%%%%%%%%%%

There are several categories of manifold structure:

\[
\boxed{
\text{smooth}
\longrightarrow
\mathrm{PL}
\longrightarrow
\text{topological}.
}
\]

A smooth manifold has differentiable transition maps.

A PL manifold has piecewise-linear transition behavior after suitable
triangulation.

A topological manifold requires only homeomorphic local coordinates.

The relationships among these structures depend strongly on dimension.

%%%%%%%%%%%%%%%%%%%%%%%%%%%%%%%%%%%%%%%%%%%%%%%%%%%%%%%%%%%%
\subsection{Exotic Smooth Structures}
\label{subsec:exotic-smooth-structures-manifolds}
%%%%%%%%%%%%%%%%%%%%%%%%%%%%%%%%%%%%%%%%%%%%%%%%%%%%%%%%%%%%

Two manifolds may be homeomorphic without being diffeomorphic.

Thus the same topological manifold may support distinct smooth structures.

An exotic sphere is homeomorphic to

\[
S^n
\]

but not diffeomorphic to the standard smooth sphere.

Even more remarkably,

\[
\boxed{
\R^4
}
\]

admits exotic smooth structures.

This makes dimension four exceptional.

%%%%%%%%%%%%%%%%%%%%%%%%%%%%%%%%%%%%%%%%%%%%%%%%%%%%%%%%%%%%
\subsection{Riemannian Completeness}
\label{subsec:riemannian-completeness-manifolds}
%%%%%%%%%%%%%%%%%%%%%%%%%%%%%%%%%%%%%%%%%%%%%%%%%%%%%%%%%%%%

Once a smooth manifold is equipped with a Riemannian metric, one can discuss
metric and geodesic completeness.

The Hopf--Rinow Theorem relates several forms of completeness for connected
Riemannian manifolds.

In particular, geodesic completeness and metric completeness are equivalent
under the theorem's standard hypotheses.

This is geometric information added to the underlying smooth manifold.

%%%%%%%%%%%%%%%%%%%%%%%%%%%%%%%%%%%%%%%%%%%%%%%%%%%%%%%%%%%%
\subsection{Every Smooth Manifold Admits a Complete Metric}
\label{subsec:complete-riemannian-metric-existence}
%%%%%%%%%%%%%%%%%%%%%%%%%%%%%%%%%%%%%%%%%%%%%%%%%%%%%%%%%%%%

Every smooth manifold satisfying the standard manifold assumptions admits a
complete Riemannian metric.

Therefore noncompactness does not prevent the existence of complete
Riemannian geometry.

For example,

\[
\R^n
\]

with its standard Euclidean metric is complete.

%%%%%%%%%%%%%%%%%%%%%%%%%%%%%%%%%%%%%%%%%%%%%%%%%%%%%%%%%%%%
\subsection{Dimension and Local Homology}
\label{subsec:dimension-local-homology}
%%%%%%%%%%%%%%%%%%%%%%%%%%%%%%%%%%%%%%%%%%%%%%%%%%%%%%%%%%%%

Manifold dimension is a topological invariant.

A fundamental reason is that neighborhoods have local topology resembling

\[
\R^n.
\]

At a point \(p\) of an \(n\)-manifold without boundary, local homology
detects the dimension through groups of the form

\[
\boxed{
H_k(M,M\setminus\{p\}).
}
\]

These groups behave like the corresponding local homology of

\[
\R^n.
\]

%%%%%%%%%%%%%%%%%%%%%%%%%%%%%%%%%%%%%%%%%%%%%%%%%%%%%%%%%%%%
\subsection{Invariance of Dimension}
\label{subsec:invariance-dimension-manifolds}
%%%%%%%%%%%%%%%%%%%%%%%%%%%%%%%%%%%%%%%%%%%%%%%%%%%%%%%%%%%%

An open subset of

\[
\R^m
\]

cannot be homeomorphic to an open subset of

\[
\R^n
\]

when

\[
m\neq n.
\]

Consequently, a connected manifold cannot simultaneously have two different
dimensions.

This ultimately follows from Brouwer's invariance of domain and related
dimension theory.

%%%%%%%%%%%%%%%%%%%%%%%%%%%%%%%%%%%%%%%%%%%%%%%%%%%%%%%%%%%%
\subsection{Invariance of Domain}
\label{subsec:invariance-of-domain-manifolds}
%%%%%%%%%%%%%%%%%%%%%%%%%%%%%%%%%%%%%%%%%%%%%%%%%%%%%%%%%%%%

\begin{theorem}[Invariance of Domain]
Let

\[
U\subseteq\R^n
\]

be open and let

\[
f:U\to\R^n
\]

be continuous and injective.

Then

\[
\boxed{
f(U)
}
\]

is open in \(\R^n\), and

\[
f:U\to f(U)
\]

is a homeomorphism.
\end{theorem}

This theorem is a fundamental reason manifold coordinates behave
consistently.

%%%%%%%%%%%%%%%%%%%%%%%%%%%%%%%%%%%%%%%%%%%%%%%%%%%%%%%%%%%%
\subsection{Classification Problems}
\label{subsec:manifold-classification-problems}
%%%%%%%%%%%%%%%%%%%%%%%%%%%%%%%%%%%%%%%%%%%%%%%%%%%%%%%%%%%%

A central problem is to classify manifolds under different notions of
equivalence:

\[
\boxed{
\text{homeomorphism},
}
\]

\[
\boxed{
\text{PL equivalence},
}
\]

\[
\boxed{
\text{diffeomorphism}.
}
\]

The answer depends strongly on dimension.

For surfaces, classification is complete and elegant.

For \(3\)-manifolds, geometrization provides a powerful structural theory.

For \(4\)-manifolds, the theory becomes exceptionally subtle.

For higher dimensions, surgery and cobordism provide major tools.

%%%%%%%%%%%%%%%%%%%%%%%%%%%%%%%%%%%%%%%%%%%%%%%%%%%%%%%%%%%%
\subsection{Manifolds in Physics}
\label{subsec:manifolds-physics}
%%%%%%%%%%%%%%%%%%%%%%%%%%%%%%%%%%%%%%%%%%%%%%%%%%%%%%%%%%%%

Manifolds provide the natural language for many physical configuration
spaces.

In general relativity, spacetime is modeled as a \(4\)-dimensional
Lorentzian manifold.

In classical mechanics, configuration spaces are manifolds, while phase
spaces frequently arise as cotangent bundles

\[
\boxed{
T^*Q.
}
\]

Gauge theories use principal bundles and connections.

Thus manifold theory is central to modern mathematical physics.

%%%%%%%%%%%%%%%%%%%%%%%%%%%%%%%%%%%%%%%%%%%%%%%%%%%%%%%%%%%%
\subsection{Manifolds in Data Analysis}
\label{subsec:manifolds-data-analysis}
%%%%%%%%%%%%%%%%%%%%%%%%%%%%%%%%%%%%%%%%%%%%%%%%%%%%%%%%%%%%

In high-dimensional data analysis, one often assumes observations lie near a
lower-dimensional manifold embedded in a large ambient space.

Schematically,

\[
\boxed{
\text{high-dimensional observations}
\approx
\text{low-dimensional manifold}.
}
\]

This motivates techniques such as

\[
\boxed{
\text{manifold learning},
\quad
\text{dimensionality reduction},
\quad
\text{geometric data analysis}.
}
\]

The mathematical manifold and the numerical approximation to it should,
however, be carefully distinguished.

%%%%%%%%%%%%%%%%%%%%%%%%%%%%%%%%%%%%%%%%%%%%%%%%%%%%%%%%%%%%
\subsection{Manifolds in Robotics}
\label{subsec:manifolds-robotics}
%%%%%%%%%%%%%%%%%%%%%%%%%%%%%%%%%%%%%%%%%%%%%%%%%%%%%%%%%%%%

Configuration spaces of mechanical systems are often manifolds.

For example, a rotating planar joint has configuration space

\[
\boxed{
S^1.
}
\]

Two independent rotating joints have configuration space

\[
\boxed{
S^1\times S^1=T^2.
}
\]

Thus even elementary robotic systems naturally produce nontrivial
manifolds.

%%%%%%%%%%%%%%%%%%%%%%%%%%%%%%%%%%%%%%%%%%%%%%%%%%%%%%%%%%%%
\subsection{Worked Example: Product Dimension}
\label{subsec:worked-product-dimension-manifold}
%%%%%%%%%%%%%%%%%%%%%%%%%%%%%%%%%%%%%%%%%%%%%%%%%%%%%%%%%%%%

\begin{workedexamplebox}{Dimension of a Product Manifold}

Consider

\[
M=S^2\times T^3.
\]

We know

\[
\dim S^2=2
\]

and

\[
\dim T^3=3.
\]

Therefore

\[
\dim M
=
2+3.
\]

\begin{answerbox}
\[
\boxed{
\dim(S^2\times T^3)=5.
}
\]
\end{answerbox}
\end{workedexamplebox}

%%%%%%%%%%%%%%%%%%%%%%%%%%%%%%%%%%%%%%%%%%%%%%%%%%%%%%%%%%%%
\subsection{Worked Example: Tangent Bundle Dimension}
\label{subsec:worked-tangent-bundle-dimension}
%%%%%%%%%%%%%%%%%%%%%%%%%%%%%%%%%%%%%%%%%%%%%%%%%%%%%%%%%%%%

\begin{workedexamplebox}{Dimension of \(TM\)}

Suppose

\[
\dim M=7.
\]

Each tangent space

\[
T_pM
\]

has dimension \(7\).

Locally,

\[
TM
\]

looks like

\[
\R^7\times\R^7.
\]

Therefore

\[
\dim TM=14.
\]

\begin{answerbox}
\[
\boxed{
\dim TM=14.
}
\]
\end{answerbox}
\end{workedexamplebox}

%%%%%%%%%%%%%%%%%%%%%%%%%%%%%%%%%%%%%%%%%%%%%%%%%%%%%%%%%%%%
\subsection{Worked Example: Quotient Dimension}
\label{subsec:worked-quotient-dimension}
%%%%%%%%%%%%%%%%%%%%%%%%%%%%%%%%%%%%%%%%%%%%%%%%%%%%%%%%%%%%

\begin{workedexamplebox}{A Free Lie Group Action}

Suppose a \(3\)-dimensional Lie group \(G\) acts freely and properly on a
\(10\)-dimensional manifold \(M\).

Then

\[
M/G
\]

is a manifold and

\[
\dim(M/G)
=
\dim M-\dim G.
\]

Therefore

\[
\dim(M/G)
=
10-3.
\]

\begin{answerbox}
\[
\boxed{
\dim(M/G)=7.
}
\]
\end{answerbox}
\end{workedexamplebox}

%%%%%%%%%%%%%%%%%%%%%%%%%%%%%%%%%%%%%%%%%%%%%%%%%%%%%%%%%%%%
\subsection{Common Errors}
\label{subsec:manifolds-common-errors}
%%%%%%%%%%%%%%%%%%%%%%%%%%%%%%%%%%%%%%%%%%%%%%%%%%%%%%%%%%%%

\subsubsection{Confusing Manifold Dimension with Ambient Dimension}

Although

\[
S^2\subseteq\R^3,
\]

the sphere has dimension

\[
\boxed{
2,
}
\]

not \(3\).

%%%%%%%%%%%%%%%%%%%%%%%%%%%%%%%%%%%%%%%%%%%%%%%%%%%%%%%%%%%%
\subsubsection{Assuming Locally Euclidean Means Globally Euclidean}

The sphere and torus are locally Euclidean but not globally homeomorphic to
Euclidean space.

%%%%%%%%%%%%%%%%%%%%%%%%%%%%%%%%%%%%%%%%%%%%%%%%%%%%%%%%%%%%
\subsubsection{Forgetting Hausdorff and Countability Conditions}

Local Euclidean structure alone is not the standard definition of a
manifold.

Hausdorff and second-countability assumptions exclude important
pathologies.

%%%%%%%%%%%%%%%%%%%%%%%%%%%%%%%%%%%%%%%%%%%%%%%%%%%%%%%%%%%%
\subsubsection{Confusing a Chart with an Atlas}

A chart describes one coordinate neighborhood.

An atlas is a collection of charts covering the manifold.

%%%%%%%%%%%%%%%%%%%%%%%%%%%%%%%%%%%%%%%%%%%%%%%%%%%%%%%%%%%%
\subsubsection{Assuming Coordinates Are Globally Defined}

Most manifolds cannot be covered by a single coordinate chart.

Coordinates are usually local.

%%%%%%%%%%%%%%%%%%%%%%%%%%%%%%%%%%%%%%%%%%%%%%%%%%%%%%%%%%%%
\subsubsection{Confusing Smooth Structure with Riemannian Metric}

A smooth structure determines which functions are differentiable.

A Riemannian metric additionally determines lengths and angles.

%%%%%%%%%%%%%%%%%%%%%%%%%%%%%%%%%%%%%%%%%%%%%%%%%%%%%%%%%%%%
\subsubsection{Assuming Every Topological Manifold Has a Unique Smooth Structure}

Smooth structures need not be unique.

Some topological manifolds admit inequivalent smooth structures.

%%%%%%%%%%%%%%%%%%%%%%%%%%%%%%%%%%%%%%%%%%%%%%%%%%%%%%%%%%%%
\subsubsection{Assuming Every Topological Manifold Is Smoothable}

Not every topological manifold admits a smooth structure.

Smoothability is a genuine additional condition.

%%%%%%%%%%%%%%%%%%%%%%%%%%%%%%%%%%%%%%%%%%%%%%%%%%%%%%%%%%%%
\subsubsection{Assuming Every Quotient Is a Manifold}

Quotient spaces can have singularities or non-Hausdorff behavior.

Conditions such as free and proper group actions are important.

%%%%%%%%%%%%%%%%%%%%%%%%%%%%%%%%%%%%%%%%%%%%%%%%%%%%%%%%%%%%
\subsubsection{Confusing Tangent Space and Tangent Bundle}

The tangent space

\[
T_pM
\]

belongs to one point.

The tangent bundle

\[
TM
\]

contains tangent spaces over every point.

%%%%%%%%%%%%%%%%%%%%%%%%%%%%%%%%%%%%%%%%%%%%%%%%%%%%%%%%%%%%
\subsubsection{Confusing Vector Fields with Individual Tangent Vectors}

A tangent vector belongs to one tangent space.

A vector field assigns tangent vectors smoothly across the manifold.

%%%%%%%%%%%%%%%%%%%%%%%%%%%%%%%%%%%%%%%%%%%%%%%%%%%%%%%%%%%%
\subsubsection{Assuming Every Vector Bundle Is Trivial}

Vector bundles may twist globally.

The M\"obius bundle is the simplest example.

%%%%%%%%%%%%%%%%%%%%%%%%%%%%%%%%%%%%%%%%%%%%%%%%%%%%%%%%%%%%
\subsubsection{Assuming Every Manifold Is Orientable}

The M\"obius strip and \(\mathbb{RP}^2\) are nonorientable.

%%%%%%%%%%%%%%%%%%%%%%%%%%%%%%%%%%%%%%%%%%%%%%%%%%%%%%%%%%%%
\subsubsection{Confusing Closed Manifold with Closed Subset}

In manifold terminology,

\[
\boxed{
\text{closed manifold}
=
\text{compact manifold without boundary}.
}
\]

%%%%%%%%%%%%%%%%%%%%%%%%%%%%%%%%%%%%%%%%%%%%%%%%%%%%%%%%%%%%
\subsubsection{Assuming Compact Means Boundaryless}

The closed disk

\[
D^n
\]

is compact but has boundary

\[
S^{n-1}.
\]

%%%%%%%%%%%%%%%%%%%%%%%%%%%%%%%%%%%%%%%%%%%%%%%%%%%%%%%%%%%%
\subsubsection{Assuming Noncompact Means Incomplete}

A noncompact manifold can admit a complete Riemannian metric.

For example,

\[
\R^n
\]

with its Euclidean metric is complete.

%%%%%%%%%%%%%%%%%%%%%%%%%%%%%%%%%%%%%%%%%%%%%%%%%%%%%%%%%%%%
\subsubsection{Assuming Every Topological Manifold Is Triangulable}

Triangulation theory for topological manifolds is dimension dependent.

Smooth manifolds are triangulable, but arbitrary topological manifolds need
not be.

%%%%%%%%%%%%%%%%%%%%%%%%%%%%%%%%%%%%%%%%%%%%%%%%%%%%%%%%%%%%
\subsection{Applications and Connections}
\label{subsec:manifolds-applications}
%%%%%%%%%%%%%%%%%%%%%%%%%%%%%%%%%%%%%%%%%%%%%%%%%%%%%%%%%%%%

\begin{applicationbox}
\textbf{Topology.}
Manifolds are among the central spaces studied by algebraic, differential,
and geometric topology.

\medskip

\textbf{Differential Geometry.}
A smooth manifold becomes a Riemannian manifold after introducing a
Riemannian metric.

\medskip

\textbf{Algebraic Topology.}
Fundamental groups, homology, cohomology, characteristic classes, and
Poincare duality reveal global manifold structure.

\medskip

\textbf{Lie Theory.}
Lie groups are manifolds whose smooth and algebraic structures interact.

\medskip

\textbf{Symplectic Geometry.}
Cotangent bundles naturally carry symplectic structures and form phase
spaces in mechanics.

\medskip

\textbf{Complex Geometry.}
Complex manifolds use local coordinates modeled on \(\mathbb{C}^n\) with
holomorphic transition functions.

\medskip

\textbf{Physics.}
Spacetime, configuration spaces, phase spaces, and gauge fields are modeled
using manifolds and bundles.

\medskip

\textbf{Robotics.}
Configuration spaces of joints and rigid motions naturally form manifolds.

\medskip

\textbf{Data Science.}
Manifold-learning methods attempt to identify lower-dimensional structure
inside high-dimensional datasets.

\medskip

\textbf{Optimization.}
Optimization can be performed on manifolds when variables satisfy nonlinear
geometric constraints.
\end{applicationbox}

%%%%%%%%%%%%%%%%%%%%%%%%%%%%%%%%%%%%%%%%%%%%%%%%%%%%%%%%%%%%
\subsection{Exercises}
\label{subsec:manifolds-exercises}
%%%%%%%%%%%%%%%%%%%%%%%%%%%%%%%%%%%%%%%%%%%%%%%%%%%%%%%%%%%%

\begin{exercise}[1]
Define a topological \(n\)-manifold.
\end{exercise}

\begin{exercise}[1]
What does locally Euclidean mean?
\end{exercise}

\begin{exercise}[1]
Why is the Hausdorff condition included in the standard manifold
definition?
\end{exercise}

\begin{exercise}[1]
What is second countability?
\end{exercise}

\begin{exercise}[1]
Define a coordinate chart.
\end{exercise}

\begin{exercise}[1]
Define an atlas.
\end{exercise}

\begin{exercise}[1]
Define a transition map.
\end{exercise}

\begin{exercise}[1]
What is a smooth atlas?
\end{exercise}

\begin{exercise}[1]
What is a smooth structure?
\end{exercise}

\begin{exercise}[1]
What is a manifold with boundary?
\end{exercise}

\begin{exercise}[1]
Define a tangent space.
\end{exercise}

\begin{exercise}[1]
Define the tangent bundle.
\end{exercise}

\begin{exercise}[1]
Define the cotangent space.
\end{exercise}

\begin{exercise}[1]
Define a vector field.
\end{exercise}

\begin{exercise}[1]
Define a fiber bundle.
\end{exercise}

\begin{exercise}[1]
Define a section of a fiber bundle.
\end{exercise}

\begin{exercise}[1]
Define a vector bundle.
\end{exercise}

\begin{exercise}[1]
Define a parallelizable manifold.
\end{exercise}

\begin{exercise}[1]
What is a closed manifold?
\end{exercise}

\begin{exercise}[2]
Determine the dimension of

\[
S^4\times T^3.
\]
\end{exercise}

\begin{exercise}[2]
Determine the dimension of

\[
S^2\times S^5\times S^1.
\]
\end{exercise}

\begin{exercise}[2]
If

\[
\dim M=8,
\]

determine

\[
\dim TM.
\]
\end{exercise}

\begin{exercise}[2]
If

\[
\dim M=6,
\]

determine

\[
\dim T^*M.
\]
\end{exercise}

\begin{exercise}[2]
Explain why \(S^1\) is a \(1\)-manifold.
\end{exercise}

\begin{exercise}[2]
Explain why \(S^2\) is a \(2\)-manifold.
\end{exercise}

\begin{exercise}[2]
Explain why an open subset of an \(n\)-manifold is itself an
\(n\)-manifold.
\end{exercise}

\begin{exercise}[2]
Show that

\[
\partial D^4=S^3.
\]
\end{exercise}

\begin{exercise}[2]
Find

\[
\partial([0,1]\times S^1).
\]
\end{exercise}

\begin{exercise}[2]
Explain why

\[
T^2=S^1\times S^1
\]

is a \(2\)-manifold.
\end{exercise}

\begin{exercise}[2]
Explain why

\[
\mathrm{GL}(2,\R)
\]

is a \(4\)-dimensional manifold.
\end{exercise}

\begin{exercise}[2]
Determine the real dimension of

\[
\mathbb{CP}^5.
\]
\end{exercise}

\begin{exercise}[2]
Explain why

\[
\mathbb{RP}^3
\]

is orientable.
\end{exercise}

\begin{exercise}[2]
Explain why

\[
\mathbb{RP}^2
\]

is nonorientable.
\end{exercise}

\begin{exercise}[2]
Give an example of a nontrivial vector bundle.
\end{exercise}

\begin{exercise}[2]
Explain why a vector field is a section of the tangent bundle.
\end{exercise}

\begin{exercise}[2]
Show that the \(n\)-torus is parallelizable.
\end{exercise}

\begin{exercise}[2]
Explain why every connected manifold is path connected.
\end{exercise}

\begin{exercise}[2]
Determine the universal cover of \(T^2\).
\end{exercise}

\begin{exercise}[3]
Construct explicit coordinate charts covering \(S^1\).
\end{exercise}

\begin{exercise}[3]
Construct stereographic coordinate charts covering \(S^2\).
\end{exercise}

\begin{exercise}[3]
Compute the transition map between stereographic charts on \(S^1\) or
\(S^2\).
\end{exercise}

\begin{exercise}[3]
Prove that

\[
\dim(M\times N)
=
\dim M+\dim N.
\]
\end{exercise}

\begin{exercise}[3]
Use the Regular Value Theorem to prove that

\[
S^n
\]

is an \(n\)-dimensional smooth submanifold of \(\R^{n+1}\).
\end{exercise}

\begin{exercise}[3]
Show that

\[
\mathrm{SL}(n,\R)
=
\{A\in M_n(\R):\det A=1\}
\]

is a smooth manifold and determine its dimension.
\end{exercise}

\begin{exercise}[3]
Show that

\[
T^n
\cong
\R^n/\Z^n.
\]
\end{exercise}

\begin{exercise}[3]
Explain why a free and proper action is important when constructing quotient
manifolds.
\end{exercise}

\begin{exercise}[3]
Prove that the boundary of an \(n\)-manifold with boundary is an
\((n-1)\)-manifold.
\end{exercise}

\begin{exercise}[3]
Explain geometrically why

\[
\partial(\partial M)=\varnothing.
\]
\end{exercise}

\begin{exercise}[3]
Construct the tangent bundle of \(S^1\) explicitly and show that it is
trivial.
\end{exercise}

\begin{exercise}[3]
Show that

\[
TS^1\cong S^1\times\R.
\]
\end{exercise}

\begin{exercise}[3]
Explain why the M\"obius line bundle over \(S^1\) is not trivial.
\end{exercise}

\begin{exercise}[3]
Prove that every Lie group is parallelizable.
\end{exercise}

\begin{exercise}[3]
Explain how a partition of unity can be used to construct a Riemannian
metric.
\end{exercise}

\begin{exercise}[3]
Show that a nowhere-vanishing top-dimensional form determines an orientation.
\end{exercise}

\begin{exercise}[3]
Explain why the existence of a global frame implies that a vector bundle is
trivial.
\end{exercise}

\begin{exercise}[3]
Show that

\[
H_n(M;\Z)\cong\Z
\]

for a familiar closed connected orientable \(n\)-manifold.
\end{exercise}

\begin{exercise}[3]
Verify Poincar\'e duality for the sphere \(S^n\).
\end{exercise}

\begin{exercise}[3]
Verify Poincar\'e duality for \(T^2\).
\end{exercise}

\begin{exercise}[4]
Study a proof of Brouwer's Invariance of Domain.
\end{exercise}

\begin{exercise}[4]
Research the proof that smooth manifolds admit partitions of unity.
\end{exercise}

\begin{exercise}[4]
Study a proof that every smooth manifold admits a Riemannian metric.
\end{exercise}

\begin{exercise}[4]
Research the proof that every smooth manifold admits a complete Riemannian
metric.
\end{exercise}

\begin{exercise}[4]
Study the Whitney Embedding Theorem.
\end{exercise}

\begin{exercise}[4]
Research the relationship between smooth, PL, and topological structures in
different dimensions.
\end{exercise}

\begin{exercise}[4]
Investigate examples of topological manifolds that are not smoothable.
\end{exercise}

\begin{exercise}[4]
Research exotic smooth structures on spheres.
\end{exercise}

\begin{exercise}[4]
Research exotic smooth structures on \(\R^4\).
\end{exercise}

\begin{exercise}[4]
Study the tangent bundle of \(S^2\) and explain why it is nontrivial.
\end{exercise}

\begin{exercise}[4]
Research why only

\[
S^0,\quad S^1,\quad S^3,\quad S^7
\]

are parallelizable among spheres.
\end{exercise}

\begin{exercise}[4]
Study the first Stiefel--Whitney class and explain how it detects
orientability.
\end{exercise}

\begin{exercise}[4]
Research the Euler class and its relationship with the Poincare--Hopf
Theorem.
\end{exercise}

\begin{exercise}[4]
Study Chern classes of complex vector bundles.
\end{exercise}

\begin{exercise}[4]
Research the frame bundle of a smooth manifold.
\end{exercise}

\begin{exercise}[4]
Study principal bundles and explain their role in gauge theory.
\end{exercise}

\begin{exercise}[4]
Research the proof of Poincare duality for closed oriented manifolds.
\end{exercise}

%%%%%%%%%%%%%%%%%%%%%%%%%%%%%%%%%%%%%%%%%%%%%%%%%%%%%%%%%%%%
\subsection{Conceptual Summary}
\label{subsec:manifolds-summary}
%%%%%%%%%%%%%%%%%%%%%%%%%%%%%%%%%%%%%%%%%%%%%%%%%%%%%%%%%%%%

A manifold is a space that is locally modeled on Euclidean space.

A topological \(n\)-manifold is locally homeomorphic to

\[
\boxed{
\R^n.
}
\]

Coordinate charts

\[
(U,\varphi)
\]

provide local numerical descriptions.

Collections of compatible charts form atlases.

Smooth transition maps give a smooth structure.

Thus

\[
\boxed{
\text{topological manifold}
+
\text{smooth structure}
=
\text{smooth manifold}.
}
\]

The manifold dimension is intrinsic and need not equal the dimension of an
ambient Euclidean space.

Manifolds can be constructed through

\[
\boxed{
\text{products},
\quad
\text{quotients},
\quad
\text{submanifolds},
\quad
\text{gluing}.
}
\]

Each point \(p\) has a tangent space

\[
\boxed{
T_pM,
}
\]

and these spaces assemble into the tangent bundle

\[
\boxed{
TM.
}
\]

Their duals form

\[
\boxed{
T^*M.
}
\]

Vector fields are sections of \(TM\), while differential \(1\)-forms are
sections of \(T^*M\).

Partitions of unity provide a fundamental local-to-global mechanism.

They allow local constructions to be combined into global structures,
including Riemannian metrics.

Fiber bundles generalize the local-product principle:

\[
\boxed{
\pi^{-1}(U)\cong U\times F.
}
\]

Vector bundles, tangent bundles, cotangent bundles, covering spaces, and
principal bundles all fit into this framework.

Characteristic classes detect global twisting of bundles.

Closed oriented manifolds possess a fundamental class and satisfy
Poincare duality:

\[
\boxed{
H^k(M)\cong H_{n-k}(M).
}
\]

Manifolds therefore provide a meeting point for

\[
\boxed{
\text{topology},
\quad
\text{geometry},
\quad
\text{analysis},
\quad
\text{algebra},
\quad
\text{physics}.
}
\]

%%%%%%%%%%%%%%%%%%%%%%%%%%%%%%%%%%%%%%%%%%%%%%%%%%%%%%%%%%%%
\subsection{Section Connections}
\label{subsec:manifolds-connections}
%%%%%%%%%%%%%%%%%%%%%%%%%%%%%%%%%%%%%%%%%%%%%%%%%%%%%%%%%%%%

\chapterconnection{
Manifolds provide the central class of spaces connecting the major branches
of modern geometry and topology. Point-Set Topology supplies the Hausdorff,
countability, compactness, connectedness, and quotient-space foundations.
Algebraic Topology studies manifolds through fundamental groups, homology,
cohomology, fundamental classes, and Poincare duality. Differential Topology
develops smooth maps, embeddings, regular values, transversality, vector
fields, Morse theory, and cobordism. Geometric Topology studies manifold
classification, decompositions, surgery, and dimension-dependent phenomena.
Differential Geometry adds metrics, connections, curvature, and geodesics.
Complex Geometry replaces real coordinates with holomorphic complex
coordinates. Symplectic Geometry places canonical structures on spaces such
as cotangent bundles. Lie Theory studies manifolds carrying compatible group
operations. Fiber bundles and characteristic classes provide major bridges
between topology, geometry, and mathematical physics. These ideas prepare
the reader for the specialized study of Knot Theory and Low-Dimensional
Topology.
}

%%%%%%%%%%%%%%%%%%%%%%%%%%%%%%%%%%%%%%%%%%%%%%%%%%%%%%%%%%%%
% SECTION SOURCE TRACKING
%%%%%%%%%%%%%%%%%%%%%%%%%%%%%%%%%%%%%%%%%%%%%%%%%%%%%%%%%%%%

%------------------------------------------------------------
% 4.5 Manifolds
%------------------------------------------------------------
%
% Source basis:
%   - Standard topological manifold, smooth manifold,
%     differential topology, and introductory bundle theory.
%
% Used for:
%   - manifold motivation;
%   - topological manifolds;
%   - Hausdorff and second-countability conditions;
%   - local Euclidean structure;
%   - dimension;
%   - coordinate charts;
%   - coordinate functions;
%   - atlases;
%   - transition maps;
%   - smooth compatibility;
%   - smooth atlases and smooth structures;
%   - maximal atlases;
%   - Euclidean spaces;
%   - circles, spheres, and tori;
%   - stereographic charts;
%   - product manifolds;
%   - quotient manifolds;
%   - projective spaces;
%   - manifolds with boundary;
%   - submanifolds;
%   - regular level sets;
%   - smooth maps and diffeomorphisms;
%   - tangent and cotangent spaces;
%   - tangent and cotangent bundles;
%   - vector fields;
%   - differential forms;
%   - immersions, submersions, and embeddings;
%   - Whitney embedding;
%   - orientability;
%   - orientation forms;
%   - partitions of unity;
%   - Riemannian metric existence;
%   - Lie groups;
%   - fiber bundles;
%   - vector bundles;
%   - sections;
%   - trivial bundles;
%   - parallelizable manifolds;
%   - covering spaces;
%   - principal bundles;
%   - frame bundles;
%   - characteristic classes;
%   - Stiefel--Whitney classes;
%   - Chern classes;
%   - Euler class;
%   - compact and closed manifolds;
%   - connectedness;
%   - universal covers;
%   - homology of manifolds;
%   - Poincare duality;
%   - CW-complex connections;
%   - triangulations;
%   - topological, PL, and smooth categories;
%   - exotic smooth structures;
%   - completeness;
%   - invariance of dimension;
%   - invariance of domain;
%   - manifold classification;
%   - applications in physics, robotics, data analysis,
%     and optimization.
%
% Notes:
%   - Point-set foundations are developed canonically in
%     Point-Set Topology.
%   - Fundamental groups, homology, cohomology, and
%     Poincare duality belong canonically to Algebraic
%     Topology.
%   - Regular values, transversality, Morse theory, and
%     cobordism belong canonically to Differential Topology.
%   - Surface and 3-manifold classification belongs to
%     Geometric Topology.
%   - Metrics, connections, curvature, and completeness
%     belong canonically to Differential Geometry.
%   - Fiber bundles and characteristic classes are
%     introduced structurally here and may be developed
%     more deeply in later topology and geometry material.
%   - Complex manifolds and complex vector bundles connect
%     to Complex Geometry.
%   - Cotangent bundles connect directly to Symplectic
%     Geometry.
%
%%%%%%%%%%%%%%%%%%%%%%%%%%%%%%%%%%%%%%%%%%%%%%%%%%%%%%%%%%%%